% !TEX root = ../main.tex
% ---------------------------------------------------------------------
%  B8 -- Funktionale Sicherheit, TwinSAFE (Skript Kap. 10; AS 9)
% ---------------------------------------------------------------------
\vlsec{cB8}{B8}{Funktionale Sicherheit und TwinSAFE}
\label{sec:safety}

\SW{Funktionale Sicherheit} Die Steuerung schützt Personen und Umwelt,
nicht nur den Betrieb. Oft sicherer, kontrolliert zu stoppen als einfach
die Energie zu kappen \ra Not-Halt, Schutztür, Zweihand in die SPS.
\begin{ztab}{@{}lX@{}}
Norm & Einstufung \\ \midrule
IEC 61508 & E/E/PE-Systeme allgemein: \T{SIL 1--4} (4 $=$ höchste) \\
ISO 13849-1 & Maschinen: \T{PL a--e}, Kategorien B, 1--4 \\
\end{ztab}
\begin{formel}[Grundprinzip]
\T{Redundanz} (1oo2): zwei Kanäle/Controller rechnen und vergleichen
sich gegenseitig. Abweichung \ra \T{Fail-Safe}: sicherer Zustand
(abschalten statt fehlerhaft weiterlaufen).
\T{Diagnosedeckungsgrad}: wie viele Fehler erkannt werden.
\T{Common Cause Failure}: eine Ursache (Überspannung) darf nicht beide
Kanäle gleichzeitig ausschalten \ra getrennte Versorgung.
\end{formel}

\subsection{TwinSAFE-Klemmen}
\begin{ztab}{@{}lX@{}}
EL6900 & TwinSAFE-Logic: „zweite SPS“ für Safety, Programm läuft auf
  \emph{zwei} Controllern, Ergebnisse verglichen; bis 128 Verbindungen;
  SIL 3 / PL e \\
EL1904 & 4 sichere Eingänge, zwei Controller mit eigener Versorgung,
  Testpulse erkennen Drahtbruch (nur mit NC-Kontakten, \emph{Sensortest
  active}) \\
EL2904 & 4 sichere Ausgänge, 24\,V, 0.5\,A je Kanal, doppelt ausgeführt \\
\end{ztab}
Gelb markiert, Bezeichnung \R{x9xx}. Sicherheitsprotokoll eingebettet im
normalen EtherCAT-Telegramm \ra gleicher Bus wie Standard-Klemmen.

\begin{rezept}[Safety-Projekt (AS 9)]
\begin{enumerate}
\item E/A scannen: EL6900, EL1904, EL2904 sichtbar.
\item SAFETY \ra neues TwinSAFE-Projekt, Target \emph{Hardware Safety
  PLC}, Physical Device $=$ EL6900.
\item TwinSAFE-Gruppe, Alias Devices EL1904/EL2904 hinzufügen und mit
  der Klemme verbinden.
\item Logik in der \R{.sal}-Datei \T{grafisch} aus validierten
  \R{safe}-Bausteinen (nicht ST); Ein-/Ausgänge benennen, im Variable
  Mapping verknüpfen.
\end{enumerate}
\end{rezept}
\begin{falle}[Not-Aus \ra Relais]
Normalzustand: Not-Aus nicht gedrückt, Relais \emph{an}. Ruhestromprinzip:
Drahtbruch oder Stromausfall muss wie „Not-Aus gedrückt“ wirken
(NC-Kontakt, Testpulse) \ra Relais \emph{aus}. Nach dem Lösen nicht
automatisch wieder anlaufen: Quittierung (Reset-Taster) verlangen.
\end{falle}
